\begin{summarybox}
\textbf{\textcolor{CSummary}{一句话核心}}
\textbf{横轴型双曲线：动点到两个实轴端点的连线斜率，乘积恒正，等于\(b^2/a^2=e^2-1\)。}
\[
\boxed{
\frac{x^2}{a^2}-\frac{y^2}{b^2}=1\quad\Longrightarrow\quad
k_{PA_1}k_{PA_2}=\frac{y^2}{x^2-a^2}=\frac{b^2}{a^2}=e^2-1.
}
\]

\textbf{\textcolor{CSummary}{三个方法怎么选}}
\begin{description}[style=nextline,leftmargin=2.8em,itemsep=0.25em]
 \item[\textbf{最快（考场）：}]斜率积直接写成\(y^2/(x^2-a^2)\)，拿曲线方程整体代换。
 \item[\textbf{最直观（入门）：}]作垂足\(H\)，两条斜率是“同一竖直高÷两条水平底”；相乘出现平方差。
 \item[\textbf{最规范（书写）：}]设点\(P(x_0,y_0)\)\(\to\)列斜率\(\to\)代入标准方程。
\end{description}

\textbf{\textcolor{CSummary}{使用前先核对四件事}}
\begin{itemize}[leftmargin=*,itemsep=0.28em]
 \item 定点是\textbf{两实轴顶点}，非焦点；\(P\)确在曲线上。
 \item \(P\)不是顶点，所需斜率均存在。
 \item 横轴型\(k_1k_2=b^2/a^2\)，纵轴型\(k_1k_2=a^2/b^2\)。
 \item \textbf{双曲线正、椭圆负}（此处两点均取长轴/实轴端点）。
\end{itemize}

\textbf{\textcolor{CSummary}{两项高价值迁移}}
\[
\boxed{e=\sqrt{1+k_1k_2}}\ \text{（横轴型）};\qquad
\boxed{\frac{x^2}{a^2}-\frac{y^2}{a^2m}=1}\ \text{（已知积}\ m>0\text{反求轨迹）}.
\]
若还知道斜率和\(s\)，则\(k_1,k_2\)是\(t^2-st+m=0\)的两根，
但两条连线必须\(k_1\ne k_2\)，因此\(\boxed{s^2>4m}\)。

\textbf{\textcolor{CSummary}{最小反例（防止套错）：}}
在\(\frac{x^2}{4}-\frac{y^2}{5}=1\)上，\(P=(3,\frac52)\)，
两斜率为\(\frac12\)和\(\frac52\)，乘积\(\frac54>0\)，
\textbf{并非}\(-\frac54\)。

\medskip
\textbf{\textcolor{CSummary}{新增两条直觉：用初中的圆连接双曲线}}
\begin{itemize}[leftmargin=*,itemsep=.2em]
 \item \textbf{圆幂法：}以 $A_1A_2$ 为直径作辅助圆；从 $H$ 引切线 $HT$，
 $HA_1\cdot HA_2=HT^2$，于是 $\boxed{k_1k_2=(PH/HT)^2=b^2/a^2}$。
 \item \textbf{圆周角法：}$k_1k_2=-1\iff PA_1\perp PA_2$，
 轨迹为以 $A_1A_2$ 为直径的\textbf{圆（删去两个端点）}。
\end{itemize}
\textbf{\textcolor{CSummary}{同一对端点的统一轨迹公式}}
\[
 \boxed{k_1k_2=m\iff y^2=m(x^2-a^2)}\quad(x\ne\pm a).
\]
\textbf{$m<0$：椭圆（$m=-1$ 为圆）；$m>0$：双曲线；$m=0$：$x$ 轴。}
以上轨迹全都要\textbf{排除 $A_1,A_2$}。

\textbf{\textcolor{CSummary}{六种证明视角：按使用场景选择}}
\begin{description}[style=nextline,leftmargin=2.8em,itemsep=.18em]
 \item[\textbf{考场最快（方法一）：}]斜率积 $y^2/(x^2-a^2)$，用方程整体替换。
 \item[\textbf{最低门槛（方法二）：}]两直角三角形，共同高除以水平底。
 \item[\textbf{初中衔接（方法三）：}]圆幂定理将两个水平底之积变成切线长平方。
 \item[\textbf{常规规范（方法四）：}]设点、斜率、代入方程、检查定义域。
 \item[\textbf{逆向轨迹（方法五）：}]从 $k_1k_2=m$ 倒推圆锥曲线（双曲线 $m>0$）。
 \item[\textbf{高阶巧证（方法六）：}]$k_1=(b/a)\tan(\theta/2)$，$k_2=(b/a)\cot(\theta/2)$。
\end{description}
\textbf{\textcolor{CSummary}{渐近线与严格判别式}}
若 $m>0$、$s=k_1+k_2$，有 $\boxed{s^2>4m}$。
$|s|=2\sqrt m$ 只是\textbf{趋向渐近线时的极限}，有限点不可取得。

\end{summarybox}
